% !TeX program = lualatex
% ============================================================
% chapter6.tex — ch06 唯一內容源（2026-08-09 起；LaTeX 統一 U1）
% 沿革：升格自 dist/ch06/chapter6.tex（原 make_dist.py 自 fragment 轉換的成品；
%   fragment 樹已凍結，改課文一律改本檔。拍板見 ../../KICKOFF-latex-unification.md）
%   樣式源＝../../template/calcbook.sty（M-B1 拍板模板）
% 編譯（本資料夾為 CWD；aux 進 build/、成品 PDF 移入 ../../dist/ch06/）：
%   latexmk -lualatex -auxdir=../../build/aux-ch06 chapter6.tex
% ============================================================
\documentclass[a4paper,12pt,oneside]{memoir}
\makeatletter\def\input@path{{../../template/}}\makeatother
\def\cbfontsdir{../../template/fonts/inter/}
\usepackage{calcbook}
% 圖：emitter 射 `<ch>/<stem>`（見 convert.py），故 graphicspath 指向 chapters/ 上層。
% 模板內的 {../figs/} 是 chapters/<ch>/ 為 CWD 時的路徑，dist/<ch>/ 需這一行覆寫。
\graphicspath{{../../chapters/}}
\begin{document}
\chapteropener{Chapter 6}{Integrals}
\begin{lead}
The first five chapters were about the derivative, the instantaneous rate at which one quantity changes as another does. This chapter turns to the opposite question, and to a problem far older than the derivative: how to measure the \emph{total amount} accumulated from a rate, or the \emph{area} enclosed by a curve. The Greeks computed areas of curved regions by approximating them ever more closely with polygons; the same idea, sharpened into a limit, becomes the \emph{definite integral}. We build it from the ground up as a limit of sums. Then, in the Fundamental Theorem of Calculus, we discover that this process of summation and the derivative are inverse to each other. That link between areas and slopes ties the two halves of calculus together, and this chapter is where we prove it.
\end{lead}
\parahead{By the end of this chapter, you will be able to:}
\begin{objectives}
  \item express areas and distances as limits of Riemann sums;
  \item define the definite integral and evaluate simple cases directly from the definition;
  \item state and apply both parts of the Fundamental Theorem of Calculus;
  \item use the Fundamental Theorem to evaluate definite integrals through antiderivatives;
  \item find indefinite integrals and read the integral of a rate as a net change;
  \item evaluate integrals by the substitution rule, including definite integrals and problems with symmetry.
\end{objectives}
\sechead{6.1}{Areas and Distances}
We know the area of a rectangle, a triangle, or any polygon: cut it into pieces whose areas we already have and add them up. But what is the area of a region with a \emph{curved} boundary, say the region under the parabola \(y = x^{2}\) from \(x = 0\) to \(x = 1\)? There is no elementary formula, because there is no way to tile a curved region exactly with finitely many rectangles. The way forward is the oldest idea in the subject: \emph{approximate} the region by rectangles we \emph{can} measure, make the approximation finer, and see what number the areas approach. That number, when it exists, is what we \emph{mean} by the area. The same method will, a little later, compute a distance from a record of velocities. The two problems turn out to be one.

\subsechead{The area problem}
Let \(f\) be continuous and non-negative on \([a, b]\), and let \(S\) be the region bounded above by \(y = f(x)\), below by the \(x\)-axis, and on the sides by the vertical lines \(x = a\) and \(x = b\). To approximate the area of \(S\), divide \([a, b]\) into \(n\) strips of equal width

\[ \Delta x = \frac{b - a}{n}, \]

by the points \(x_{0} = a,\ x_{1},\ x_{2},\ \dots,\ x_{n} = b\), where \(x_{i} = a + i\,\Delta x\). Over each strip \([x_{i-1}, x_{i}]\) build a rectangle whose height is the value of \(f\) at the \emph{right} endpoint, \(f(x_{i})\). The combined area of these \(n\) rectangles is the \emph{right endpoint sum}

\[ R_{n} = f(x_{1})\,\Delta x + f(x_{2})\,\Delta x + \dots + f(x_{n})\,\Delta x = \sum_{i=1}^{n} f(x_{i})\,\Delta x, \]

where \(\sum_{i=1}^{n}\) is the sigma notation for the sum of the listed terms as \(i\) runs from \(1\) to \(n\). Choosing instead the \emph{left} endpoint of each strip gives the \emph{left endpoint sum} \(L_{n} = \sum_{i=1}^{n} f(x_{i-1})\,\Delta x\). Each is a sum of rectangle areas, an approximation to the area of \(S\) that we expect to sharpen as the strips grow thinner.

\begin{workedexample}
\begin{envexample}{Example}{6.1}{}
Estimate the area under \(y = x^{2}\) from \(0\) to \(1\) using \(n = 4\) rectangles of equal width, with heights taken at the right endpoints. Then do the same with the left endpoints, and read off what the two estimates say about the true area.
\end{envexample}
\begin{envsolution}{Solution}{}{}
Here \(\Delta x = \tfrac{1}{4}\), and the division points are \(0, \tfrac14, \tfrac12, \tfrac34, 1\). With \(f(x) = x^{2}\), the right endpoints are \(\tfrac14, \tfrac12, \tfrac34, 1\), so

\[ \begin{aligned}
          R_{4} &= \tfrac{1}{4}\!\left[\left(\tfrac14\right)^{2} + \left(\tfrac12\right)^{2} + \left(\tfrac34\right)^{2} + 1^{2}\right] \\[2pt]
                &= \tfrac{1}{4}\cdot\frac{1 + 4 + 9 + 16}{16} = \frac{30}{64} = \frac{15}{32} \approx 0.469.
        \end{aligned} \]

The left endpoints are \(0, \tfrac14, \tfrac12, \tfrac34\), which drop the tallest rectangle and add a flat one of height \(0\):

\[ \begin{aligned}
          L_{4} &= \tfrac{1}{4}\!\left[0^{2} + \left(\tfrac14\right)^{2} + \left(\tfrac12\right)^{2} + \left(\tfrac34\right)^{2}\right] \\[2pt]
                &= \tfrac{1}{4}\cdot\frac{0 + 1 + 4 + 9}{16} = \frac{14}{64} = \frac{7}{32} \approx 0.219.
        \end{aligned} \]

Because \(x^{2}\) is increasing on \([0, 1]\), each right rectangle rises above the curve and each left rectangle falls below it, so the two sums bracket the true area \(A\):

\[ 0.219 \approx L_{4} < A < R_{4} \approx 0.469. \]

The direction of this bracket is a fact about the \emph{increase} of \(x^{2}\), not about which endpoint is read. On an interval where \(f\) \emph{decreases}, the roles swap: the right endpoints undershoot and the left endpoints overshoot. Where \(f\) both rises and falls, over- and under-estimation must be judged piece by piece.

With four rectangles, the left and right sums bracket the area only loosely. The bracket will tighten as \(n\) grows, and the next example takes \(n\) all the way to infinity. \qedmark
\end{envsolution}
\end{workedexample}
\begin{figureblock}
  \includegraphics[width=66.67mm]{ch06/figs/riemann-lr-x2-1}\hspace{6mm}\includegraphics[width=66.67mm]{ch06/figs/riemann-lr-x2-2}
\figcaption{Figure 6.1}{Approximating the area under \(y = x^{2}\) on \([0, 1]\) with four equal strips. Reading each rectangle’s height at the \emph{left} endpoint (left panel) keeps every rectangle below the curve, so \(L_{4} = \tfrac{7}{32}\) underestimates; reading it at the \emph{right} endpoint (right panel) lifts every rectangle above the curve, so \(R_{4} = \tfrac{15}{32}\) overestimates. The true area is trapped between.}
\end{figureblock}
To let \(n \to \infty\) we need the right endpoint sum as an explicit function of \(n\). For \(f(x) = x^{2}\) on \([0, 1]\) the width is \(\Delta x = \tfrac{1}{n}\) and the right endpoints are \(x_{i} = \tfrac{i}{n}\), so \(f(x_{i}) = \tfrac{i^{2}}{n^{2}}\) and

\[ R_{n} = \sum_{i=1}^{n} \frac{i^{2}}{n^{2}}\cdot\frac{1}{n} = \frac{1}{n^{3}} \sum_{i=1}^{n} i^{2}. \]

The remaining sum \(\sum_{i=1}^{n} i^{2}\) is a \emph{power sum}. Unlike the arithmetic sum \(\sum_{i=1}^{n} i = \tfrac{n(n+1)}{2}\), it has the less familiar closed form

\[ \sum_{i=1}^{n} i^{2} = \frac{n(n+1)(2n+1)}{6}, \]

proved by telescoping in Appendix A (Proposition A.6). With that closed form in hand, the limit is a routine computation.

\begin{figureblock}
  \begin{minipage}{\linewidth}\centering
  \includegraphics[width=57.41mm]{ch06/figs/refinement-rn-x2-1}\hspace{6mm}\includegraphics[width=57.41mm]{ch06/figs/refinement-rn-x2-2}\\[6pt]
  \includegraphics[width=57.41mm]{ch06/figs/refinement-rn-x2-3}
  \figcaption{Figure 6.2}{Refining the right endpoint sum for \(y = x^{2}\) on \([0, 1]\). With \(n = 4, 8, 16\) strips the rectangles overshoot the curve on every strip, but by less and less as \(n\) grows, so the overestimates close in on the true area from above.}
  \end{minipage}
\end{figureblock}
\begin{workedexample}
\begin{envexample}{Example}{6.2}{}
Find the exact area under \(y = x^{2}\) from \(0\) to \(1\) as the limit of the right endpoint sums.
\end{envexample}
\begin{envsolution}{Solution}{}{}
Substitute the closed form for \(\sum i^{2}\) into \(R_{n}\) and simplify, keeping the powers of \(n\) grouped so the limit is easy to read:

\[ \begin{aligned}
          R_{n} &= \frac{1}{n^{3}} \sum_{i=1}^{n} i^{2} \\[2pt]
                &= \frac{1}{n^{3}}\cdot\frac{n(n+1)(2n+1)}{6} \\[2pt]
                &= \frac{(n+1)(2n+1)}{6 n^{2}} \\[2pt]
                &= \frac{1}{6}\left(1 + \frac{1}{n}\right)\!\left(2 + \frac{1}{n}\right).
        \end{aligned} \]

As \(n \to \infty\), both \(\tfrac{1}{n}\) terms vanish, so

\[ A = \lim_{n \to \infty} R_{n} = \frac{1}{6}\,(1)(2) = \frac{1}{3}. \]

The exact area is \(\tfrac{1}{3}\), comfortably inside the bracket \(0.219 < A < 0.469\) from Example 6.1. The same limit results from the left endpoint sums \(L_{n}\): they differ from \(R_{n}\) by only the single end rectangle \(\tfrac{1}{n}\bigl(f(1) - f(0)\bigr) = \tfrac{1}{n}\), which itself tends to \(0\). Approaching the area from below or from above, we arrive at the one number. \qedmark
\end{envsolution}
\end{workedexample}
That the left and right sums converge to a single value is exactly what lets us \emph{define} the area to be that value.

\begin{envdefinition}{Definition}{6.1}{Area under a curve}
Let \(f\) be continuous and non-negative on \([a, b]\). The \emph{area} of the region under the graph of \(f\) from \(a\) to \(b\) is

\[ A = \lim_{n \to \infty} \sum_{i=1}^{n} f(x_{i})\,\Delta x, \qquad \Delta x = \frac{b - a}{n},\quad x_{i} = a + i\,\Delta x, \]

provided this limit exists. For a continuous \(f\) it does, and the left endpoint sums, the right endpoint sums, and in fact the sums built from \emph{any} choice of sample point in each strip all approach the same value.
\end{envdefinition}
\begin{envcaution}{Caution}{}{}
The clause “provided this limit exists” is not a formality. It rests on a genuine theorem. Continuity of \(f\) guarantees that the approximating sums converge to one limit, whatever sample points are chosen.

We rely on this guarantee here so that the definition gives one unambiguous value. The next section states the theorem precisely, says what its proof would need, and takes it on credit. It is the one result of this chapter we do not prove.

For a badly discontinuous \(f\), sums built from different choices of sample point really can disagree. Driving every sample to where the function misbehaves gives one value, and steering clear of those places gives another. No single limit exists, and hence no area is defined. (For the equal-width left and right sums themselves the gap is only \(\bigl(f(b) - f(a)\bigr)\,\Delta x\), which closes regardless. Continuity is needed to handle the freedom in the sample points, not the choice of side.)
\end{envcaution}
The method is old. In the third century BCE, Archimedes found the area of a parabolic segment by trapping it between inscribed and circumscribed straight-sided figures, then bringing the two bounds together. That is the same bracket, computed by hand, that Examples 6.1 and 6.2 tightened into a limit. What the seventeenth century added was not the idea but the machinery to carry the limit out in general, which is what this chapter does.

\subsechead{The distance problem}
Suppose we know the \emph{speed} of a moving object at every instant, but not its position, and we want the \emph{total distance} it travels over a time interval. Speed is never negative, because it measures how fast the object moves, whichever way it is heading. We take it to vary continuously. If it were \emph{constant}, the distance would be just speed times elapsed time. It is not constant, but over a \emph{short} stretch of time it is nearly so, and on that stretch \(\text{speed} \times \text{time}\) is a good estimate of the distance covered. Add these estimates over the whole interval and the total distance is approximated by a sum of exactly the form \(\sum v(t_{i})\,\Delta t\), a Riemann sum again, now with speed in place of \(f\) and elapsed time in place of width.

\begin{workedexample}
\begin{envexample}{Example}{6.3}{}
A runner starts from rest and speeds up, her speed recorded every second for five seconds:

\[ \begin{array}{c|cccccc} t\ (\text{s}) & 0 & 1 & 2 & 3 & 4 & 5 \\ \hline v\ (\text{m/s}) & 0 & 4 & 7 & 9 & 10 & 10 \end{array} \]

Assuming her speed climbs continuously and never falls back over this interval, estimate the distance she covers.
\end{envexample}
\begin{envsolution}{Solution}{}{}
Take \(\Delta t = 1\) second. Using the speed at the \emph{left} end of each one-second interval (the readings at \(t = 0, 1, 2, 3, 4\)) gives the underestimate

\[ L = 1\cdot(0 + 4 + 7 + 9 + 10) = 30 \text{ m}, \]

while the speed at the \emph{right} end of each interval (the readings at \(t = 1, 2, 3, 4, 5\)) gives the overestimate

\[ R = 1\cdot(4 + 7 + 9 + 10 + 10) = 40 \text{ m}. \]

Because the speed never falls back, on each one-second stretch it stays no lower than the left reading and no higher than the right, so the left sum is a lower estimate of the distance and the right sum an upper one: the true distance lies between them, \(30 \text{ m} \le \text{distance} \le 40 \text{ m}\), about \(35\) metres. The readings alone fix only the sampled values; the bracket rests on the modelling assumption that the speed climbs without falling back between them. Readings taken every half-second, then every tenth of a second, would squeeze the two estimates together. That is the same refinement that Example 6.2 carried to the limit. \qedmark
\end{envsolution}
\end{workedexample}
\begin{figureblock}
  \includegraphics[width=82.66mm]{ch06/figs/velocity-distance-steps}
\figcaption{Figure 6.3}{The runner’s recorded speeds at one-second marks (Example 6.3), with the right endpoint rectangles whose areas sum to the overestimate \(R\). Since distance accumulates speed times time, the distance travelled is the area under the speed–time graph; the dashed line only suggests the unrecorded speed between readings.}
\end{figureblock}
Set the two computations side by side. The area under a curve is \(\lim_{n\to\infty}\sum f(x_{i})\,\Delta x\); the distance from a speed is \(\lim_{n\to\infty}\sum v(t_{i})\,\Delta t\). They are the \emph{same} limit of the \emph{same} kind of sum, differing only in what the quantities are called. Indeed, for a non-negative speed the distance is exactly the area under the speed-versus-time graph. Whenever a quantity accumulates from a varying rate, the same construction returns: chop the interval, multiply rate by width, sum, and refine. The next section names this common limit the \emph{definite integral} and introduces its symbol. It also states the theorem that guarantees the limit for every continuous function.


\sechead{6.2}{The Definite Integral}
In Section 6.1 two different problems led to the \emph{same} limit of the \emph{same} kind of sum: the area under a curve, and the distance travelled at a varying velocity. That is no coincidence. A construction that keeps reappearing is worth naming, and worth a symbol of its own. Here we name the common limit the \emph{definite integral} and introduce its symbol. We also remove the assumption that \(f\) is non-negative. For a function that dips below the axis the same limit still makes perfect sense; it simply measures area with a sign, and that signed quantity is what the integral records.

Fix a function \(f\) on \([a, b]\). Exactly as before, cut \([a, b]\) into \(n\) strips of equal width \(\Delta x = (b - a)/n\) by points \(a = x_{0} < x_{1} < \dots < x_{n} = b\); but now, in each strip \([x_{i-1}, x_{i}]\), allow the height to be read at \emph{any} point \(x_{i}^{*}\) of that strip, not just an endpoint. The resulting sum

\[ \sum_{i=1}^{n} f(x_{i}^{*})\,\Delta x \]

is called a \emph{Riemann sum} for \(f\) on \([a, b]\), after Bernhard Riemann. The left, right, and midpoint sums are the special cases \(x_{i}^{*} = x_{i-1}\), \(x_{i}^{*} = x_{i}\), and \(x_{i}^{*} = \tfrac{1}{2}(x_{i-1} + x_{i})\). The definite integral is what these sums approach.

\begin{envdefinition}{Definition}{6.2}{Definite integral}
Let \(f\) be defined on \([a, b]\). The \emph{definite integral of \(f\) from \(a\) to \(b\)} is

\[ \int_{a}^{b} f(x)\,dx = \lim_{n \to \infty} \sum_{i=1}^{n} f(x_{i}^{*})\,\Delta x, \qquad \Delta x = \frac{b - a}{n}, \]

provided this limit exists and takes the \emph{same} value for every choice of the sample points \(x_{i}^{*}\). When it does, \(f\) is called \emph{integrable} on \([a, b]\).
\end{envdefinition}
The symbol \(\int\), an elongated \(S\) for \emph{summa}, was Leibniz’s; it still shows the structure of the Riemann sum, even though the limit itself no longer refers to individual rectangles. In \(\displaystyle\int_{a}^{b} f(x)\,dx\), the numbers \(a\) and \(b\) are the \emph{limits of integration} (lower and upper), \(f(x)\) is the \emph{integrand}, and the \(dx\) marks the variable with respect to which we sum (what is left of the width \(\Delta x\)). The integral is a \emph{number}, not a function of \(x\); the \(x\) inside is a placeholder, so

\[ \int_{a}^{b} f(x)\,dx = \int_{a}^{b} f(t)\,dt = \int_{a}^{b} f(u)\,du. \]

Leibniz settled on the elongated \(S\) in 1675, writing the integral as a continuous analogue of a sum of infinitely many infinitely thin rectangles. He based the notation on infinitesimals, an idea that was hard to make precise. The notation is still in use because it translates directly into the limit of Definition 6.2.

To use the integral freely we extend it to the cases \(a = b\) and \(a > b\) by convention, so that the limits of integration may sit in either order:

\[ \int_{a}^{a} f(x)\,dx = 0, \qquad \int_{b}^{a} f(x)\,dx = -\int_{a}^{b} f(x)\,dx. \]

Reversing the limits reverses the sign, since running the strips from \(b\) back to \(a\) turns each width \(\Delta x\) into its negative.

\subsechead{When does the limit exist?}
Definition 6.2 is a strong requirement: \emph{every} choice of sample points must give the same number. For a badly behaved function this fails. Recall the caution of §6.1: there, some choices of sample point fell where \(f\) misbehaves and others steered clear of those places, so different choices gave different answers. For the functions we actually meet, continuity prevents this.

\begin{envtheorem}{Theorem}{6.1}{Integrability of continuous functions}
If \(f\) is continuous on \([a, b]\), then \(f\) is integrable on \([a, b]\): the Riemann sums approach a single limit, the same for every way of slicing \([a, b]\) into strips whose widths shrink to zero, and every choice of sample points within them. The same holds if \(f\) is bounded on \([a, b]\) and continuous there except at finitely many points.
\end{envtheorem}
\begin{envcaution}{Caution}{}{}
Theorem 6.1 is the one result of this chapter we do not prove. Its proof turns on a property of a continuous function on a \emph{closed} interval that is subtler than continuity at each point. That property is \emph{uniform} continuity, a single rate of continuity good across the whole interval, and it forces the gap between the largest and smallest Riemann sums to shrink to zero as the strips narrow. We take the theorem on credit and use what it gives us: for a continuous integrand the limit is always there, and the sample points may be chosen for convenience. Everything else in this chapter is proved in full.
\end{envcaution}
\subsechead{Net signed area}
When \(f\) can be negative, a strip where \(f(x_{i}^{*}) < 0\) contributes a \emph{negative} term \(f(x_{i}^{*})\,\Delta x\) to the Riemann sum. So the integral counts area above the \(x\)-axis as positive and area below it as negative:

\[ \int_{a}^{b} f(x)\,dx = \bigl(\text{area above the axis}\bigr) - \bigl(\text{area below the axis}\bigr), \]

the \emph{net signed area} between the graph and the axis. For a non-negative \(f\) nothing lies below the axis, and this reduces to the area of Definition 6.1. The definite integral extends that idea to functions that change sign.

\begin{workedexample}
\begin{envexample}{Example}{6.4}{}
Evaluate \(\displaystyle\int_{0}^{1} x^{2}\,dx\) directly from the definition.
\end{envexample}
\begin{envsolution}{Solution}{}{}
The integrand \(x^{2}\) is continuous, so Theorem 6.1 guarantees that the limit exists and that it does not depend on the choice of sample points. We may therefore use the right-endpoint sums already computed in §6.1 (Example 6.2), where they were found to converge to \(\tfrac{1}{3}\). Hence

\[ \int_{0}^{1} x^{2}\,dx = \frac{1}{3}. \]

What §6.1 called “the area under \(y = x^{2}\)” is now, with the axis label supplied, the value of a definite integral. Because \(x^{2} \ge 0\), the signed area is the plain area. \qedmark
\end{envsolution}
\end{workedexample}
\begin{workedexample}
\begin{envexample}{Example}{6.5}{}
Evaluate \(\displaystyle\int_{0}^{2} (x - 1)\,dx\) using net signed area.
\end{envexample}
\begin{envsolution}{Solution}{}{}
The graph of \(y = x - 1\) is a straight line crossing the axis at \(x = 1\). Over \([0, 1]\) it lies \emph{below} the axis, bounding a triangle with legs \(1\) and \(1\), of area \(\tfrac{1}{2}\); over \([1, 2]\) it lies \emph{above} the axis, bounding a congruent triangle of area \(\tfrac{1}{2}\). The integral is the area above minus the area below:

\[ \int_{0}^{2} (x - 1)\,dx = \tfrac{1}{2} - \tfrac{1}{2} = 0. \]

The two signed pieces cancel exactly. Note the contrast with the \emph{total} area enclosed between the line and the axis, which is \(\tfrac{1}{2} + \tfrac{1}{2} = 1\): the integral reports the net, not the total. \qedmark
\end{envsolution}
\end{workedexample}
\begin{figureblock}
  \includegraphics[width=79.9mm]{ch06/figs/signed-area-line}
\figcaption{Figure 6.4}{The integral \(\int_{0}^{2} (x - 1)\,dx\) as net signed area. The line crosses the axis at \(x = 1\); the triangle below the axis (marked \(-\), area \(\tfrac{1}{2}\)) cancels the congruent triangle above it (marked \(+\), area \(\tfrac{1}{2}\)), so the integral is \(0\) even though the total area enclosed is \(1\).}
\end{figureblock}
\begin{workedexample}
\begin{envexample}{Example}{6.6}{}
Evaluate \(\displaystyle\int_{-1}^{1} \sqrt{1 - x^{2}}\,dx\) by interpreting it as an area.
\end{envexample}
\begin{envsolution}{Solution}{}{}
The graph of \(y = \sqrt{1 - x^{2}}\) is the upper half of the unit circle: squaring gives \(y^{2} = 1 - x^{2}\), that is \(x^{2} + y^{2} = 1\), and the requirement \(y \ge 0\) keeps only the top. Over \([-1, 1]\) the region between this graph and the axis is therefore a semicircular disc of radius \(1\), of area \(\tfrac{1}{2}\pi(1)^{2} = \tfrac{\pi}{2}\). Since the integrand is non-negative, the definite integral is exactly that area:

\[ \int_{-1}^{1} \sqrt{1 - x^{2}}\,dx = \frac{\pi}{2}. \]

When an integrand’s graph bounds a region whose area we already know, we can read the integral straight off the geometry, with no limit to compute. \qedmark
\end{envsolution}
\end{workedexample}
\begin{figureblock}
  \includegraphics[width=53.8mm]{ch06/figs/semicircle-area}
\figcaption{Figure 6.5}{Reading \(\int_{-1}^{1} \sqrt{1 - x^{2}}\,dx\) as an area (Example 6.6). The graph \(y = \sqrt{1 - x^{2}}\) is the upper half of the unit circle, so the shaded region is a semicircle of radius \(1\); its area \(\tfrac{\pi}{2}\) is the value of the integral.}
\end{figureblock}
\subsechead{Properties of the integral}
Because the integral is a limit of sums, the arithmetic of sums passes to it in the limit. The following properties hold for continuous \(f\) and \(g\); each is the limit of the corresponding statement about Riemann sums.

\begin{envtheorem}{Theorem}{6.2}{Properties of the definite integral}
Let \(f\) and \(g\) be continuous on \([a, b]\) with \(a \le b\); let \(k\) be a constant; and let \(c\) satisfy \(a \le c \le b\).

\begin{steps}
  \item \emph{Constant.} \(\displaystyle\int_{a}^{b} k\,dx = k\,(b - a)\).
  \item \emph{Linearity.} \(\displaystyle\int_{a}^{b}\bigl[f(x) + g(x)\bigr]dx = \int_{a}^{b} f(x)\,dx + \int_{a}^{b} g(x)\,dx\) and \(\displaystyle\int_{a}^{b} k\,f(x)\,dx = k\int_{a}^{b} f(x)\,dx\).
  \item \emph{Additivity over intervals.} \(\displaystyle\int_{a}^{c} f(x)\,dx + \int_{c}^{b} f(x)\,dx = \int_{a}^{b} f(x)\,dx\).
  \item \emph{Comparison.} If \(f(x) \ge 0\) on \([a, b]\), then \(\displaystyle\int_{a}^{b} f(x)\,dx \ge 0\). More generally, if \(f(x) \le g(x)\) on \([a, b]\), then \(\displaystyle\int_{a}^{b} f(x)\,dx \le \int_{a}^{b} g(x)\,dx\); and if \(m \le f(x) \le M\) on \([a, b]\), then \(\displaystyle m\,(b - a) \le \int_{a}^{b} f(x)\,dx \le M\,(b - a)\).
\end{steps}
\end{envtheorem}
\emph{Constant and scalar rules.} Each is a single line: \(\sum k\,\Delta x = k\,(b - a)\), since the \(n\) equal widths total \(n\,\Delta x = b - a\), and \(\sum k\,f(x_{i}^{*})\,\Delta x = k \sum f(x_{i}^{*})\,\Delta x\).

\emph{Linearity.} Every Riemann sum splits termwise, \(\sum \bigl[f(x_{i}^{*}) + g(x_{i}^{*})\bigr]\Delta x = \sum f(x_{i}^{*})\,\Delta x + \sum g(x_{i}^{*})\,\Delta x\), and the limit of a sum is the sum of the limits (limit laws, Theorem 1.2).

\emph{Comparison.} If \(f \ge 0\) then every term \(f(x_{i}^{*})\,\Delta x \ge 0\), so each Riemann sum is \(\ge 0\), and a limit of non-negative numbers is non-negative. Applying this to \(g - f\), which is \(\ge 0\) exactly when \(f \le g\), and using linearity gives \(\int f \le \int g\). Squeezing \(f\) between the constants \(m\) and \(M\) then yields the two-sided bound through the constant rule.

\emph{Additivity.} This is the geometric fact that the signed area over \([a, b]\) is the signed area over \([a, c]\) together with that over \([c, b]\): the region divides at the vertical line \(x = c\). Being continuous on \([a, b]\), \(f\) is continuous on \([a, c]\) and \([c, b]\) as well, so both pieces are themselves integrals. A sum-based proof slices \([a, b]\) with a division point at \(c\), where the Riemann sum splits outright into a part over \([a, c]\) and a part over \([c, b]\), and Theorem 6.1 guarantees the same limit whichever way we slice.

The sign convention \(\int_{b}^{a} = -\int_{a}^{b}\) then extends the identity to a \(c\) lying outside \([a, b]\), as long as \(f\) is continuous on the whole interval that \(a\), \(b\), and \(c\) span.

\begin{workedexample}
\begin{envexample}{Example}{6.7}{}
Given that \(\displaystyle\int_{0}^{1} x^{2}\,dx = \tfrac{1}{3}\) (Example 6.4), and without evaluating any new limit, find \(\displaystyle\int_{0}^{1} \bigl(3x^{2} - 4\bigr)\,dx\).
\end{envexample}
\begin{envsolution}{Solution}{}{}
The properties turn one known integral into many. Linearity (Theorem 6.2, part 2) splits the integral and pulls out the constant factor, and the constant rule (part 1) evaluates the remaining piece directly:

\[ \begin{aligned}
          \int_{0}^{1} \bigl(3x^{2} - 4\bigr)\,dx &= 3\int_{0}^{1} x^{2}\,dx - \int_{0}^{1} 4\,dx \\[2pt]
                                                   &= 3\cdot\tfrac{1}{3} - 4\,(1 - 0) = 1 - 4 = -3.
        \end{aligned} \]

No new Riemann sum was evaluated: the single value \(\int_{0}^{1} x^{2} = \tfrac{1}{3}\), run through linearity, settles it. This is how the integral is used in practice. Most integrals are reached not from the definition but by reducing them to ones already known. \qedmark
\end{envsolution}
\end{workedexample}
\begin{workedexample}
\begin{envexample}{Example}{6.8}{}
Show that \(\displaystyle 3 \le \int_{1}^{4} \sqrt{x}\,dx \le 6\) without evaluating the integral.
\end{envexample}
\begin{envsolution}{Solution}{}{}
Since \(\sqrt{x}\) is continuous on \([1, 4]\), it is integrable there by Theorem 6.1, so the properties apply. On \([1, 4]\) the integrand is increasing, running from \(\sqrt{1} = 1\) up to \(\sqrt{4} = 2\), so \(1 \le \sqrt{x} \le 2\) throughout, and the bounding property (Theorem 6.2, part 4) with \(m = 1\), \(M = 2\), and \(b - a = 3\) gives

\[ 1 \cdot 3 \le \int_{1}^{4} \sqrt{x}\,dx \le 2 \cdot 3, \qquad\text{that is}\qquad 3 \le \int_{1}^{4} \sqrt{x}\,dx \le 6. \]

A single comparison, with no limit computed, traps the integral between \(3\) and \(6\). The exact value, \(\tfrac{14}{3} \approx 4.67\), which §6.3 will make easy to find, indeed sits inside. \qedmark
\end{envsolution}
\end{workedexample}
\begin{figureblock}
  \includegraphics[width=79.9mm]{ch06/figs/minmax-sqrt}
\figcaption{Figure 6.6}{Bounding \(\int_{1}^{4} \sqrt{x}\,dx\) without evaluating it (Example 6.8). On \([1, 4]\) the integrand runs from its minimum \(m = 1\) to its maximum \(M = 2\); the inscribed rectangle (area \(1 \cdot 3 = 3\)) and the circumscribed rectangle (area \(2 \cdot 3 = 6\)) trap the integral, so \(3 \le \int_{1}^{4} \sqrt{x}\,dx \le 6\).}
\end{figureblock}
We now have the definite integral on a firm footing: a number defined as a limit of Riemann sums, guaranteed to exist whenever the integrand is continuous, measuring net signed area, and obeying the linearity, additivity, and comparison rules that make it possible to reason about it. Yet every value we have produced has come the hard way, either from a limit of sums, as in Example 6.4, or from a picture, as in Example 6.5. Only a few integrals can be evaluated either way. The next section removes this obstacle: the Fundamental Theorem of Calculus connects the definite integral to the derivative, and makes the great majority of integrals easy to evaluate.


\sechead{6.3}{The Fundamental Theorem of Calculus}
Two sections of work have built the definite integral into a well-defined number, but left it hard to compute. Every value so far came from a limit of Riemann sums or from a picture. The \emph{Fundamental Theorem of Calculus} removes that difficulty, and it is the reason the subject is one method rather than two. It comes in two parts. The first says that integrating a continuous function and then differentiating recovers the function: differentiation undoes integration. The second turns that fact into a tool, evaluating a definite integral from any single antiderivative of the integrand, with no sums at all. Both parts rest on a single function: the one that measures accumulated area.

The reverse of differentiation has a name of its own.

\begin{envdefinition}{Definition}{6.3}{Antiderivative}
A function \(F\) is an \emph{antiderivative} of \(f\) on an open interval \(I\) if \(F'(x) = f(x)\) for every \(x\) in \(I\).
\end{envdefinition}
For example, \(\tfrac{1}{3}x^{3}\) is an antiderivative of \(x^{2}\), since its derivative is \(x^{2}\); so is \(\tfrac{1}{3}x^{3} + 5\). Reading each differentiation rule of Chapters 2–4 backwards produces an antiderivative, and the Fundamental Theorem is about to make antiderivatives the key to every definite integral.

\subsechead{The accumulation function}
Fix a continuous function \(f\) on \([a, b]\) and, for each \(x\) in \([a, b]\), record the signed area accumulated from \(a\) up to \(x\):

\[ g(x) = \int_{a}^{x} f(t)\,dt. \]

The variable of integration is written \(t\) to keep it distinct from the upper limit \(x\), which is now the input of a new function \(g\). As \(x\) advances, \(g(x)\) accumulates more area. How fast? Nudging \(x\) rightward by a small \(h\) appends a thin strip under the graph, of width \(h\) and height about \(f(x)\), so \(g\) gains roughly \(f(x)\,h\). Its rate of increase is the current height \(f(x)\). That guess is exactly right, and it is the first half of the theorem.

\begin{figureblock}
  \includegraphics[width=85.19mm]{ch06/figs/accumulation-sliver}
\figcaption{Figure 6.7}{The accumulation function \(g(x) = \int_{a}^{x} f(t)\,dt\) is the shaded area up to \(x\). Advancing \(x\) by a small \(h\) adds the thin strip over \([x, x + h]\), of width \(h\) and height about \(f(x)\), so \(g\) grows by roughly \(f(x)\,h\). Its rate of change is the current height \(f(x)\). The strip is widened for visibility.}
\end{figureblock}
\begin{envtheorem}{Theorem}{6.3}{Fundamental Theorem of Calculus, Part 1}
If \(f\) is continuous on \([a, b]\), then the function \(g\) defined by

\[ g(x) = \int_{a}^{x} f(t)\,dt, \qquad a \le x \le b, \]

is continuous on \([a, b]\) and differentiable on \((a, b)\), and

\[ g'(x) = f(x). \]

In Leibniz notation, \(\displaystyle\frac{d}{dx}\int_{a}^{x} f(t)\,dt = f(x)\): the accumulation function \(g\) is an antiderivative of \(f\) on \((a, b)\).
\end{envtheorem}
\begin{envproof}{Proof}{}{}
Fix \(x\) in \((a, b)\) and take \(h \ne 0\) small enough that \(x + h\) lies in \([a, b]\). By additivity of the integral (Theorem 6.2),

\[ \begin{aligned}
        g(x + h) - g(x) &= \int_{a}^{x+h} f(t)\,dt - \int_{a}^{x} f(t)\,dt \\[2pt]
                        &= \int_{x}^{x+h} f(t)\,dt.
      \end{aligned} \]

Consider the difference quotient \(\dfrac{g(x + h) - g(x)}{h} = \dfrac{1}{h}\displaystyle\int_{x}^{x+h} f(t)\,dt\). We claim it lies between the least and greatest values of \(f\) on the closed interval joining \(x\) and \(x + h\). For \(h > 0\), the comparison bound (Theorem 6.2, part 4) applied on \([x, x+h]\) gives, with \(m_{h}\) and \(M_{h}\) the minimum and maximum of \(f\) there,

\[ m_{h}\,h \le \int_{x}^{x+h} f(t)\,dt \le M_{h}\,h, \]

and dividing by \(h > 0\) yields \(m_{h} \le \dfrac{g(x+h) - g(x)}{h} \le M_{h}\). For \(h < 0\) the same bracketing holds on the interval \([x + h, x]\): with \(m_{h}, M_{h}\) now its minimum and maximum, \(m_{h}\lvert h\rvert \le \int_{x+h}^{x} f \le M_{h}\lvert h\rvert\); since \(\int_{x}^{x+h} f = -\int_{x+h}^{x} f\) and \(h = -\lvert h\rvert\), the quotient \(\dfrac{1}{h}\int_{x}^{x+h} f\) again lies between \(m_{h}\) and \(M_{h}\). Now by the Extreme Value Theorem (Theorem 4.9(a)), \(m_{h} = f(u_{h})\) and \(M_{h} = f(v_{h})\) for some points \(u_{h}, v_{h}\) in the interval. As \(h \to 0\), that interval collapses onto \(x\), forcing \(u_{h} \to x\) and \(v_{h} \to x\), so by the continuity of \(f\),

\[ m_{h} = f(u_{h}) \to f(x) \qquad\text{and}\qquad M_{h} = f(v_{h}) \to f(x). \]

The difference quotient is squeezed between two quantities both tending to \(f(x)\), so by the Squeeze Theorem (Theorem 1.3),

\[ g'(x) = \lim_{h \to 0} \frac{g(x + h) - g(x)}{h} = f(x). \]

It remains to see that \(g\) is continuous on the whole closed interval, endpoints included. Since \(f\) is continuous on \([a, b]\) it is bounded there, say \(\lvert f\rvert \le M\) (Theorem 4.9(a)). Then for \(x\) and \(x + h\) in \([a, b]\) the comparison bound gives \(\bigl\lvert g(x+h) - g(x)\bigr\rvert = \bigl\lvert\int_{x}^{x+h} f\bigr\rvert \le M\,\lvert h\rvert\), which tends to \(0\) with \(h\). So \(g\) is continuous at every point of \([a, b]\). \qedmark
\end{envproof}
\begin{figureblock}
  \includegraphics[width=85.64mm]{ch06/figs/ftc-trap}
\figcaption{Figure 6.8}{The trapping step in the proof of Part 1. Over \([x, x + h]\) the added area \(\int_{x}^{x+h} f\) lies between the inscribed rectangle of height \(m_{h} = \min f\) and the circumscribed rectangle of height \(M_{h} = \max f\), attained at points \(u_{h}, v_{h}\) in the interval. Dividing by \(h\) brackets the difference quotient between \(m_{h}\) and \(M_{h}\); as \(h \to 0\) both tend to \(f(x)\).}
\end{figureblock}
\begin{envcaution}{Caution}{}{}
Two points. First, the letter under the integral sign is only a placeholder, so \(\int_{a}^{x} f(t)\,dt\) and \(\int_{a}^{x} f(s)\,ds\) are the same function of \(x\). Writing \(\int_{a}^{x} f(x)\,dx\), with \(x\) doing double duty as both limit and dummy, is a genuine error. Second, Part 1 requires nothing of \(f\) beyond continuity, so it differentiates accumulation functions we cannot evaluate: \(g(x) = \int_{0}^{x} e^{-t^{2}}\,dt\) has \(g'(x) = e^{-x^{2}}\) exactly, even though \(e^{-t^{2}}\) has no antiderivative expressible in elementary functions. The theorem \emph{constructs} an antiderivative where no formula for one exists.
\end{envcaution}
\begin{workedexample}
\begin{envexample}{Example}{6.9}{}
Differentiate: (a) \(\displaystyle g(x) = \int_{1}^{x} \sqrt{1 + t^{2}}\,dt\); and (b) \(\displaystyle h(x) = \int_{0}^{x^{3}} \cos t\,dt\).
\end{envexample}
\begin{envsolution}{Solution}{}{}
(a) The integrand \(\sqrt{1 + t^{2}}\) is continuous, so Part 1 applies directly:

\[ g'(x) = \sqrt{1 + x^{2}}. \]

(b) The upper limit is now \(x^{3}\), not \(x\), so \(h\) is a composition. Write \(H(u) = \int_{0}^{u} \cos t\,dt\), so that \(h(x) = H(x^{3})\). Part 1 gives \(H'(u) = \cos u\), and the chain rule (Theorem 3.3) supplies the outer factor:

\[ h'(x) = H'(x^{3})\cdot \frac{d}{dx}\bigl(x^{3}\bigr) = \cos(x^{3})\cdot 3x^{2} = 3x^{2}\cos(x^{3}). \]

A variable limit of integration is differentiated by Part 1; a \emph{function} in the limit calls for the chain rule on top. \qedmark
\end{envsolution}
\end{workedexample}
\begin{workedexample}
\begin{envexample}{Example}{6.10}{}
Differentiate (a) \(\displaystyle g(x) = \int_{x}^{1} \sqrt{1 + t^{2}}\,dt\); and (b) \(\displaystyle h(x) = \int_{x^{2}}^{x^{3}} \cos t\,dt\).
\end{envexample}
\begin{envsolution}{Solution}{}{}
(a) The variable now sits in the \emph{lower} limit. Reversing the limits flips the sign, \(g(x) = -\int_{1}^{x} \sqrt{1 + t^{2}}\,dt\), and Part 1 differentiates the reversed integral:

\[ g'(x) = -\sqrt{1 + x^{2}}. \]

The minus sign matters: advancing the \emph{lower} limit reduces the accumulated area, so the rate of change is the negative of the integrand.

(b) Both limits are functions of \(x\). Split at any constant — \(0\) will do — into two variable-upper-limit pieces, \(h(x) = \int_{0}^{x^{3}} \cos t\,dt - \int_{0}^{x^{2}} \cos t\,dt\), and apply Part 1 with the chain rule to each end:

\[ h'(x) = \cos(x^{3})\cdot 3x^{2} - \cos(x^{2})\cdot 2x = 3x^{2}\cos(x^{3}) - 2x\cos(x^{2}). \]

Each end contributes the integrand evaluated at that limit, times the limit’s own derivative; the lower end enters with a minus. \qedmark
\end{envsolution}
\end{workedexample}
\subsechead{Evaluating integrals}
Part 1 built one antiderivative of \(f\), the accumulation function. Part 2 does the reverse: it evaluates a definite integral from \emph{any} antiderivative whatsoever, and most antiderivatives are far easier to write down than \(\int_{a}^{x} f\).

\begin{envtheorem}{Theorem}{6.4}{Fundamental Theorem of Calculus, Part 2}
If \(f\) is continuous on \([a, b]\) and \(F\) is \emph{any} antiderivative of \(f\) on an open interval containing \([a, b]\), then

\[ \int_{a}^{b} f(x)\,dx = F(b) - F(a). \]
\end{envtheorem}
\begin{envproof}{Proof}{}{}
Let \(g(x) = \int_{a}^{x} f(t)\,dt\). By Part 1, \(g\) is an antiderivative of \(f\) on \((a, b)\) and is continuous on \([a, b]\); the given \(F\), being differentiable on an open interval containing \([a, b]\), is continuous there too (Theorem 2.1). Their difference \(F - g\) is therefore continuous on \([a, b]\), and on \((a, b)\) its derivative is \((F - g)' = f - f = 0\). A function continuous on a closed interval whose derivative vanishes throughout the interior is constant (Corollary 4.4), so \(F(x) = g(x) + C\) for some constant \(C\) and all \(x\) in \([a, b]\). Evaluating at the endpoints, the constant cancels:

\[ \begin{aligned}
        F(b) - F(a) &= \bigl[g(b) + C\bigr] - \bigl[g(a) + C\bigr] \\[2pt]
                    &= g(b) - g(a).
      \end{aligned} \]

Finally \(g(b) = \int_{a}^{b} f(t)\,dt\) and \(g(a) = \int_{a}^{a} f(t)\,dt = 0\), so \(F(b) - F(a) = \int_{a}^{b} f(t)\,dt\), which is the claim. \qedmark
\end{envproof}
The difference \(F(b) - F(a)\) is written \(\bigl[F(x)\bigr]_{a}^{b}\) (or \(F(x)\big|_{a}^{b}\)) when the antiderivative is displayed. Any added constant is pointless here because it cancels in the subtraction. We therefore take the simplest antiderivative available.

\begin{envstrategy}{Strategy}{6.1}{Evaluating a definite integral}
To evaluate \(\displaystyle\int_{a}^{b} f(x)\,dx\) by the Fundamental Theorem:

\begin{steps}
  \item Find an antiderivative \(F\) of \(f\) — a function with \(F' = f\) — by reading the derivative rules of Chapters 2–4 backwards.
  \item Evaluate at the upper and lower limits and subtract: \(\displaystyle\int_{a}^{b} f(x)\,dx = \bigl[F(x)\bigr]_{a}^{b} = F(b) - F(a)\).
  \item Any antiderivative works, since the additive constant cancels — so choose the simplest.
\end{steps}
\end{envstrategy}
\begin{workedexample}
\begin{envexample}{Example}{6.11}{}
Evaluate \(\displaystyle\int_{0}^{1} x^{2}\,dx\) and \(\displaystyle\int_{1}^{4} \sqrt{x}\,dx\) by the Fundamental Theorem, and compare with the earlier sections.
\end{envexample}
\begin{envsolution}{Solution}{}{}
An antiderivative of \(x^{2}\) is \(\tfrac{1}{3}x^{3}\), so

\[ \int_{0}^{1} x^{2}\,dx = \left[\frac{x^{3}}{3}\right]_{0}^{1} = \frac{1}{3} - 0 = \frac{1}{3}. \]

The whole limit-of-Riemann-sums computation of §6.1–6.2 reduces to one line. For the second, write \(\sqrt{x} = x^{1/2}\), whose antiderivative is \(\tfrac{x^{3/2}}{3/2} = \tfrac{2}{3}x^{3/2}\):

\[ \begin{aligned}
          \int_{1}^{4} \sqrt{x}\,dx &= \left[\frac{2}{3}x^{3/2}\right]_{1}^{4}
                                     = \frac{2}{3}\bigl(4^{3/2} - 1^{3/2}\bigr) \\[2pt]
                                    &= \frac{2}{3}(8 - 1) = \frac{14}{3}.
        \end{aligned} \]

In Example 6.8 the bounding property could only confine this integral to the range \([3, 6]\); the Fundamental Theorem gives the exact value \(\tfrac{14}{3} \approx 4.67\), which indeed sits there. \qedmark
\end{envsolution}
\end{workedexample}
\begin{workedexample}
\begin{envexample}{Example}{6.12}{}
Evaluate \(\displaystyle\int_{0}^{\pi} \sin x\,dx\).
\end{envexample}
\begin{envsolution}{Solution}{}{}
Since \(\dfrac{d}{dx}(-\cos x) = \sin x\), an antiderivative of \(\sin x\) is \(-\cos x\), and

\[ \begin{aligned}
          \int_{0}^{\pi} \sin x\,dx &= \bigl[-\cos x\bigr]_{0}^{\pi}
                                     = (-\cos \pi) - (-\cos 0) \\[2pt]
                                    &= -(-1) - (-1) = 2.
        \end{aligned} \]

The answer is positive, as it must be: over \([0, \pi]\) the sine curve rises in a single hump above the axis, so its signed area is an ordinary positive area. \qedmark
\end{envsolution}
\end{workedexample}
This inverse link is the discovery, made independently by Newton in the 1660s and Leibniz in the 1670s, that founded calculus as one subject rather than two. The area problem and the tangent problem had been studied for centuries in isolation. In fact they are inverse to each other. The derivative of an accumulated area is the height, and so an area is recovered from an antiderivative. Recognizing this connection turned a collection of clever area computations into a method.

The two parts are one statement seen from opposite ends. Part 1 differentiates an integral and returns the integrand; Part 2 integrates using an antiderivative and returns a difference of its values. Together they say that differentiation and integration are inverse processes, the central fact of the subject. What remains is to make antiderivatives easy to produce and to read. The next section sets up a notation for the whole family of antiderivatives of a function and reads the Fundamental Theorem as a statement about accumulated change.


\sechead{6.4}{Indefinite Integrals and the Net Change Theorem}
The Fundamental Theorem reduces the evaluation of a definite integral to finding an antiderivative, so we now need a notation and a working table for antiderivatives. This section provides both, and then rereads the theorem as a statement about accumulated change, the form most useful in practice.

\begin{envdefinition}{Definition}{6.4}{Indefinite integral}
The \emph{indefinite integral} of \(f\), written \(\int f(x)\,dx\), denotes the family of all antiderivatives of \(f\). The statement

\[ \int f(x)\,dx = F(x) + C \]

means that \(F' = f\), where \(C\) is an arbitrary constant, the \emph{constant of integration}. As \(C\) ranges over all real numbers, \(F + C\) ranges over every antiderivative of \(f\) on the interval in question.
\end{envdefinition}
\begin{envcaution}{Caution}{}{}
Keep the two integrals apart. The \emph{definite} integral \(\int_{a}^{b} f(x)\,dx\) is a \emph{number}; the \emph{indefinite} integral \(\int f(x)\,dx\) is a \emph{family of functions}. They share a symbol because the Fundamental Theorem links them: the definite integral is evaluated \emph{through} an antiderivative. But a number and a family of functions are not the same kind of object. The limits of integration \(a\) and \(b\) are what distinguish the two on the page.
\end{envcaution}
\subsechead{A table of antiderivatives}
Every differentiation rule from Chapters 2 through 4, read backwards, is an antiderivative formula. The most useful ones (each carrying an implied \(+C\)) are worth collecting:

\[ \begin{aligned}
    &\int x^{n}\,dx = \frac{x^{n+1}}{n+1} \ \ (n \ne -1), &\qquad &\int \frac{1}{x}\,dx = \ln\lvert x\rvert, \\[4pt]
    &\int e^{x}\,dx = e^{x}, &\qquad &\int \cos x\,dx = \sin x, \\[4pt]
    &\int \sin x\,dx = -\cos x, &\qquad &\int \sec^{2} x\,dx = \tan x, \\[4pt]
    &\int \frac{1}{1 + x^{2}}\,dx = \arctan x, &\qquad &\int \frac{1}{\sqrt{1 - x^{2}}}\,dx = \arcsin x.
  \end{aligned} \]

The trigonometric and inverse-trigonometric entries are the Chapter 3 derivative rules read backwards. For instance, \(\int \sec^{2} x\,dx = \tan x\) because \((\tan x)' = \sec^{2} x\), and \(\arctan\) and \(\arcsin\) likewise reverse the derivatives found there. Each identity holds on any interval where both the integrand and the stated antiderivative are defined: the \(\sec^{2}\) formula where \(\cos x \ne 0\), the \(\arcsin\) formula on \(-1 < x < 1\), and — for a negative or fractional exponent — the power rule on an interval avoiding the points where \(x^{n}\) is undefined. The power rule covers every exponent except \(n = -1\), where the formula \(\tfrac{x^{n+1}}{n+1}\) would divide by zero. That missing case is filled by the logarithm. The reason an absolute value appears there is worth seeing.

For \(x > 0\), \(\lvert x\rvert = x\) and \(\dfrac{d}{dx}\ln x = \dfrac{1}{x}\), as established in Chapter 4 (Theorem 4.14). For \(x < 0\), \(\lvert x\rvert = -x\), and the chain rule gives

\[ \frac{d}{dx}\ln(-x) = \frac{1}{-x}\cdot(-1) = \frac{1}{x}. \]

So on either side of the origin \(\dfrac{d}{dx}\ln\lvert x\rvert = \dfrac{1}{x}\), which is why

\[ \int \frac{1}{x}\,dx = \ln\lvert x\rvert + C. \]

The absolute value extends the antiderivative to negative \(x\) as well as positive. One caution: \(1/x\) is undefined at \(x = 0\), and its domain falls into the two separate intervals \((-\infty, 0)\) and \((0, \infty)\); the formula is valid on any interval that avoids the origin, but \(C\) may take different values on the two sides.

\begin{workedexample}
\begin{envexample}{Example}{6.13}{}
Find (a) \(\displaystyle\int \bigl(x^{3} - 6x\bigr)\,dx\) and (b) \(\displaystyle\int \left(2e^{x} - \frac{5}{x}\right)dx\).
\end{envexample}
\begin{envsolution}{Solution}{}{}
Integrate term by term, using the table and the linearity that the derivative rules pass to antiderivatives. (a)

\[ \int \bigl(x^{3} - 6x\bigr)\,dx = \frac{x^{4}}{4} - 6\cdot\frac{x^{2}}{2} + C = \frac{x^{4}}{4} - 3x^{2} + C. \]

(b) The term \(5/x\) uses the logarithm entry:

\[ \int \left(2e^{x} - \frac{5}{x}\right)dx = 2e^{x} - 5\ln\lvert x\rvert + C. \]

A single constant \(C\) absorbs the constants from every term. We can always check an indefinite integral by differentiating the answer: \(\dfrac{d}{dx}\bigl(2e^{x} - 5\ln\lvert x\rvert\bigr) = 2e^{x} - \dfrac{5}{x}\), as required. \qedmark
\end{envsolution}
\end{workedexample}
\begin{workedexample}
\begin{envexample}{Example}{6.14}{}
Evaluate \(\displaystyle\int_{1}^{2} \frac{1}{x}\,dx\).
\end{envexample}
\begin{envsolution}{Solution}{}{}
An antiderivative of \(1/x\) is \(\ln\lvert x\rvert\), and on \([1, 2]\) the variable is positive, so \(\ln\lvert x\rvert = \ln x\). By the Fundamental Theorem,

\[ \int_{1}^{2} \frac{1}{x}\,dx = \bigl[\ln\lvert x\rvert\bigr]_{1}^{2} = \ln 2 - \ln 1 = \ln 2. \]

The area under \(y = 1/x\) from \(1\) to \(2\) is exactly \(\ln 2 \approx 0.693\). \qedmark
\end{envsolution}
\end{workedexample}
\begin{workedexample}
\begin{envexample}{Example}{6.15}{}
Evaluate \(\displaystyle\int_{0}^{1} \frac{1}{1 + x^{2}}\,dx\).
\end{envexample}
\begin{envsolution}{Solution}{}{}
The table records \(\int \dfrac{1}{1 + x^{2}}\,dx = \arctan x\), the reverse of \((\arctan x)' = \dfrac{1}{1 + x^{2}}\) established in Chapter 3. By the Fundamental Theorem,

\[ \int_{0}^{1} \frac{1}{1 + x^{2}}\,dx = \bigl[\arctan x\bigr]_{0}^{1} = \arctan 1 - \arctan 0 = \frac{\pi}{4}. \]

A rational integrand with no area to read off has produced the number \(\tfrac{\pi}{4}\), one of the standard ways of obtaining \(\pi\) from an integral. \qedmark
\end{envsolution}
\end{workedexample}
\subsechead{The Net Change Theorem}
Read the Fundamental Theorem with the integrand named as a derivative. If a quantity \(F\) changes at the rate \(F'\), then integrating that rate over \([a, b]\) recovers the total change in \(F\).

\begin{envtheorem}{Theorem}{6.5}{Net Change Theorem}
If \(F'\) is continuous on \([a, b]\), then

\[ \int_{a}^{b} F'(x)\,dx = F(b) - F(a). \]

The integral of a rate of change is the net change in the quantity.
\end{envtheorem}
This is nothing more than Part 2 of the Fundamental Theorem (Theorem 6.4) with the integrand written as \(F'\): \(F\) is an antiderivative of \(F'\), so the theorem applies verbatim. What is new is how we read it. If \(V'(t)\) is the rate at which water flows into a tank, \(\int_{t_{1}}^{t_{2}} V'(t)\,dt\) is the net volume added; if \(C'(x)\) is a marginal cost, \(\int_{x_{1}}^{x_{2}} C'(x)\,dx\) is the added cost of raising production from \(x_{1}\) to \(x_{2}\). The most important case is motion.

For an object moving along a line with velocity \(v(t)\), position is an antiderivative of velocity. The \emph{displacement} is the net change in position, and the Net Change Theorem gives it as

\[ \text{displacement} = \int_{a}^{b} v(t)\,dt. \]

This is a signed quantity: motion in the negative direction subtracts. The \emph{total distance travelled}, which never subtracts, is found by integrating the \emph{speed} \(\lvert v(t)\rvert\), as in §6.1:

\[ \text{total distance} = \int_{a}^{b} \lvert v(t)\rvert\,dt. \]

To compute the distance we split the interval at the instants where \(v\) changes sign, integrate \(v\) on each piece, and add the absolute values.

\begin{workedexample}
\begin{envexample}{Example}{6.16}{}
A particle moves along a line with velocity \(v(t) = t^{2} - 4\) (in m/s) for \(0 \le t \le 3\). Find its displacement and the total distance it travels.
\end{envexample}
\begin{envsolution}{Solution}{}{}
The displacement is the integral of the velocity over \([0, 3]\):

\[ \int_{0}^{3} \bigl(t^{2} - 4\bigr)\,dt = \left[\frac{t^{3}}{3} - 4t\right]_{0}^{3} = (9 - 12) - 0 = -3 \text{ m}. \]

The particle ends \(3\) metres behind where it started. For the total distance, note \(v(t) = t^{2} - 4 = 0\) at \(t = 2\); the velocity is negative on \([0, 2)\) (the particle moves backward), zero at \(t = 2\), and positive on \((2, 3]\). Splitting there,

\[ \begin{aligned}
          \int_{0}^{2} \bigl(t^{2} - 4\bigr)\,dt &= \left[\frac{t^{3}}{3} - 4t\right]_{0}^{2} = \frac{8}{3} - 8 = -\frac{16}{3}, \\[2pt]
          \int_{2}^{3} \bigl(t^{2} - 4\bigr)\,dt &= (9 - 12) - \left(\frac{8}{3} - 8\right) = -3 + \frac{16}{3} = \frac{7}{3}.
        \end{aligned} \]

The total distance adds the absolute values of the two signed pieces:

\[ \text{distance} = \left\lvert -\frac{16}{3}\right\rvert + \left\lvert \frac{7}{3}\right\rvert = \frac{16}{3} + \frac{7}{3} = \frac{23}{3} \approx 7.67 \text{ m}. \]

The particle covers \(\tfrac{23}{3}\) metres of ground while ending only \(3\) metres from its start: the forward leg of \(\tfrac{7}{3}\) m undoes only part of the backward leg of \(\tfrac{16}{3}\) m, leaving it \(3\) metres behind. \qedmark
\end{envsolution}
\end{workedexample}
\begin{figureblock}
  \includegraphics[width=79.9mm]{ch06/figs/velocity-signed-area}
\figcaption{Figure 6.9}{Displacement versus distance for \(v(t) = t^{2} - 4\) on \([0, 3]\) (Example 6.16). The velocity is negative on \([0, 2)\) — the region below the axis, marked \(-\), of area \(\tfrac{16}{3}\) — and positive on \((2, 3]\) — the region above, marked \(+\), of area \(\tfrac{7}{3}\). The displacement is the net \(\tfrac{7}{3} - \tfrac{16}{3} = -3\) m, while the total distance adds the magnitudes, \(\tfrac{16}{3} + \tfrac{7}{3} = \tfrac{23}{3}\) m.}
\end{figureblock}
The definite integral is now, in principle, a solved problem: find an antiderivative, subtract its values, and read the result as an area or a net change. In practice a gap remains. The table above applies directly only to functions that are visibly some derivative run backwards. An integrand with an inner function tucked inside another often has no direct entry there. The next section carries the chain rule over to the antiderivative side, and with it the single most useful technique for finding antiderivatives that the table alone cannot.


\sechead{6.5}{The Substitution Rule}
The table of §6.4 antidifferentiates a function directly only when the function is visibly a derivative written plainly. Many integrands contain an inner function inside another function, the pattern to which the chain rule applies. To integrate \(\int 2x\sqrt{1 + x^{2}}\,dx\) there is no direct table entry, yet \(2x\) is exactly the derivative of the inner expression \(1 + x^{2}\), and that is no accident. Applying the chain rule in reverse gives the method of \emph{substitution}. It is the single most useful technique of integration.

\begin{envtheorem}{Theorem}{6.6}{The Substitution Rule}
If \(u = g(x)\) is differentiable and \(f\) is continuous on the range of \(g\), then

\[ \int f\bigl(g(x)\bigr)\,g'(x)\,dx = \int f(u)\,du. \]
\end{envtheorem}
\begin{envproof}{Proof}{}{}
Let \(F\) be an antiderivative of \(f\), so \(F' = f\). By the chain rule (Theorem 3.3),

\[ \frac{d}{dx}\,F\bigl(g(x)\bigr) = F'\bigl(g(x)\bigr)\,g'(x) = f\bigl(g(x)\bigr)\,g'(x). \]

So \(F(g(x))\) is an antiderivative of \(f(g(x))\,g'(x)\), which means \(\int f(g(x))\,g'(x)\,dx = F(g(x)) + C\). On the other side, \(\int f(u)\,du = F(u) + C\), and substituting \(u = g(x)\) turns this into \(F(g(x)) + C\) as well. The two are equal. \qedmark
\end{envproof}
In practice the rule is applied through the differential notation of §5.3. Setting \(u = g(x)\) gives \(du = g'(x)\,dx\), and the rule reads as the literal replacement of \(g(x)\) by \(u\) and of \(g'(x)\,dx\) by \(du\). The art is entirely in the choice of \(u\): pick an inner function whose derivative is also present, up to a constant factor.

\begin{envstrategy}{Strategy}{6.2}{Integration by substitution}
To evaluate an integral by substitution:

\begin{steps}
  \item Choose \(u = g(x)\) — usually an inner function whose derivative \(g'(x)\) also appears in the integrand, up to a constant multiple.
  \item Compute \(du = g'(x)\,dx\) and rewrite the \emph{entire} integral in terms of \(u\); every \(x\), the \(dx\) included, must disappear. If it cannot, the choice of \(u\) was wrong.
  \item Integrate in \(u\) using the table.
  \item For an indefinite integral, back-substitute \(u = g(x)\). For a definite integral, either back-substitute, or convert the limits to \(u = g(a)\) and \(u = g(b)\) and finish in \(u\) — never leave \(x\)-limits on a \(u\)-integrand.
\end{steps}
\end{envstrategy}
\begin{workedexample}
\begin{envexample}{Example}{6.17}{}
Find \(\displaystyle\int 2x\sqrt{1 + x^{2}}\,dx\).
\end{envexample}
\begin{envsolution}{Solution}{}{}
The expression inside the square root is \(1 + x^{2}\), and its derivative \(2x\) is already present in the integrand. Put \(u = 1 + x^{2}\), so \(du = 2x\,dx\). The resulting integral is listed in the table:

\[ \int 2x\sqrt{1 + x^{2}}\,dx = \int \sqrt{u}\,du = \frac{2}{3}u^{3/2} + C = \frac{2}{3}\bigl(1 + x^{2}\bigr)^{3/2} + C. \]

Differentiating the answer returns \(2x\sqrt{1 + x^{2}}\), confirming it. \qedmark
\end{envsolution}
\end{workedexample}
\begin{workedexample}
\begin{envexample}{Example}{6.18}{}
Find \(\displaystyle\int x^{2}\bigl(x^{3} + 1\bigr)^{5}\,dx\).
\end{envexample}
\begin{envsolution}{Solution}{}{}
Take \(u = x^{3} + 1\), the inner function; then \(du = 3x^{2}\,dx\). The integrand offers \(x^{2}\,dx\), not \(3x^{2}\,dx\), but the constant is easily supplied: \(x^{2}\,dx = \tfrac{1}{3}\,du\). So

\[ \int x^{2}\bigl(x^{3} + 1\bigr)^{5}\,dx = \frac{1}{3}\int u^{5}\,du = \frac{1}{3}\cdot\frac{u^{6}}{6} + C = \frac{\bigl(x^{3} + 1\bigr)^{6}}{18} + C. \]

Only a \emph{constant} factor may be moved across the integral sign in this way; a missing factor of \(x\) could not be, and would signal a poor choice of \(u\). \qedmark
\end{envsolution}
\end{workedexample}
\begin{workedexample}
\begin{envexample}{Example}{6.19}{}
Find \(\displaystyle\int \frac{2x}{1 + x^{2}}\,dx\).
\end{envexample}
\begin{envsolution}{Solution}{}{}
The derivative of the denominator is \(2x\), exactly the numerator. That is the mark of a substitution whose result is a logarithm rather than a power. Put \(u = 1 + x^{2}\), so \(du = 2x\,dx\), and the integral becomes \(\int \tfrac{1}{u}\,du\):

\[ \int \frac{2x}{1 + x^{2}}\,dx = \int \frac{du}{u} = \ln\lvert u\rvert + C = \ln\bigl(1 + x^{2}\bigr) + C, \]

where the absolute value is dropped because \(1 + x^{2} > 0\) throughout. This is the pattern \(\int \dfrac{f'(x)}{f(x)}\,dx = \ln\lvert f(x)\rvert + C\): when the numerator is the derivative of the denominator, substitution collapses the integral to the logarithm introduced in §6.4. \qedmark
\end{envsolution}
\end{workedexample}
\subsechead{Substitution in definite integrals}
A definite integral can be substituted just as an indefinite one, provided the limits are handled with care. There are two correct options: finish the antiderivative and restore \(x\) before applying the limits, or — usually cleaner — carry the substitution through the limits themselves.

\begin{envtheorem}{Theorem}{6.7}{Substitution in a definite integral}
If \(g'\) is continuous on \([a, b]\) and \(f\) is continuous on the range of \(g\), then

\[ \int_{a}^{b} f\bigl(g(x)\bigr)\,g'(x)\,dx = \int_{g(a)}^{g(b)} f(u)\,du. \]
\end{envtheorem}
\begin{envproof}{Proof}{}{}
Let \(F\) be an antiderivative of \(f\). As in Theorem 6.6, \(F(g(x))\) is an antiderivative of \(f(g(x))\,g'(x)\), so by Part 2 of the Fundamental Theorem (Theorem 6.4),

\[ \begin{aligned}
        \int_{a}^{b} f\bigl(g(x)\bigr)\,g'(x)\,dx &= \Bigl[F\bigl(g(x)\bigr)\Bigr]_{a}^{b} = F\bigl(g(b)\bigr) - F\bigl(g(a)\bigr) \\[2pt]
                                                  &= \Bigl[F(u)\Bigr]_{g(a)}^{g(b)} = \int_{g(a)}^{g(b)} f(u)\,du.
      \end{aligned} \]

Both sides equal \(F(g(b)) - F(g(a))\); the substitution merely reads the same difference off the \(u\)-axis. \qedmark
\end{envproof}
\begin{envcaution}{Caution}{}{}
When a definite integral is converted to \(u\), the limits must be converted too. Writing \(\int_{a}^{b} \dots\,du\) with the \emph{original} \(x\)-limits still attached to a \(u\)-integrand is a common and costly error: the numbers \(a\) and \(b\) are values of \(x\), not of \(u\). Either change the limits to \(g(a)\) and \(g(b)\) and evaluate in \(u\), or return to \(x\) first and use the original limits. Never blend the two.
\end{envcaution}
\begin{workedexample}
\begin{envexample}{Example}{6.20}{}
Evaluate \(\displaystyle\int_{0}^{2} 2x\bigl(1 + x^{2}\bigr)^{3}\,dx\).
\end{envexample}
\begin{envsolution}{Solution}{}{}
Let \(u = 1 + x^{2}\), so \(du = 2x\,dx\). Convert the limits: when \(x = 0\), \(u = 1\); when \(x = 2\), \(u = 5\). The integral passes entirely into \(u\):

\[ \int_{0}^{2} 2x\bigl(1 + x^{2}\bigr)^{3}\,dx = \int_{1}^{5} u^{3}\,du = \left[\frac{u^{4}}{4}\right]_{1}^{5} = \frac{625 - 1}{4} = 156. \]

Changing the limits removes any need to return to \(x\); the answer emerges directly from the \(u\)-integral. \qedmark
\end{envsolution}
\end{workedexample}
\subsechead{Symmetry}
When the interval is symmetric about the origin, the symmetry of the integrand — if it has any — can halve the work or eliminate it entirely.

Picture the two cases. For an even function, the graph on \([-a, 0]\) is the graph on \([0, a]\) reflected across the vertical axis. Each thin strip on the left therefore has a matching strip on the right with the same height and the same sign, so the matched signed areas \emph{add}: the integral is twice the integral over \([0, a]\). An odd function’s graph on \([-a, 0]\) is the \([0, a]\) graph turned a half-turn about the origin, so each left strip is the mirror of a right strip carrying the \emph{opposite} sign; the matched contributions \emph{cancel} and the integral is zero. The substitution \(u = -x\) below turns this picture into a proof.

\begin{envtheorem}{Theorem}{6.8}{Integrals of symmetric functions}
Suppose \(f\) is continuous on \([-a, a]\).

\begin{steps}
  \item If \(f\) is \emph{even} — that is, \(f(-x) = f(x)\) — then \(\displaystyle\int_{-a}^{a} f(x)\,dx = 2\int_{0}^{a} f(x)\,dx\).
  \item If \(f\) is \emph{odd} — that is, \(f(-x) = -f(x)\) — then \(\displaystyle\int_{-a}^{a} f(x)\,dx = 0\).
\end{steps}
\end{envtheorem}
\begin{envproof}{Proof}{}{}
Split at the origin: \(\int_{-a}^{a} f = \int_{-a}^{0} f + \int_{0}^{a} f\). In the first piece substitute \(u = -x\), so \(du = -dx\), and the limits \(x = -a, 0\) become \(u = a, 0\):

\[ \int_{-a}^{0} f(x)\,dx = \int_{a}^{0} f(-u)\,(-du) = \int_{0}^{a} f(-u)\,du. \]

If \(f\) is even, \(f(-u) = f(u)\), so this equals \(\int_{0}^{a} f(u)\,du\) and the two pieces add to \(2\int_{0}^{a} f\). If \(f\) is odd, \(f(-u) = -f(u)\), so it equals \(-\int_{0}^{a} f(u)\,du\), which cancels the second piece and leaves \(0\). \qedmark
\end{envproof}
\begin{workedexample}
\begin{envexample}{Example}{6.21}{}
Evaluate \(\displaystyle\int_{-1}^{1} \bigl(x^{4} + x^{3}\bigr)\,dx\).
\end{envexample}
\begin{envsolution}{Solution}{}{}
The integrand splits into an even part and an odd part. Over the symmetric interval \([-1, 1]\), the even \(x^{4}\) doubles its integral over \([0, 1]\), while the odd \(x^{3}\) integrates to zero (Theorem 6.8):

\[ \int_{-1}^{1} \bigl(x^{4} + x^{3}\bigr)\,dx = 2\int_{0}^{1} x^{4}\,dx + 0 = 2\left[\frac{x^{5}}{5}\right]_{0}^{1} = \frac{2}{5}. \]

The odd term contributes nothing at all, so recognizing this saves the work of integrating it. \qedmark
\end{envsolution}
\end{workedexample}
Substitution is the reverse of the chain rule, and the first entry in the toolkit of integration techniques. It applies to integrands with a single inner function whose derivative is present, up to a constant factor. This is a wide class, but it does not include all integrands. Products of unrelated functions, powers of trigonometric functions, and rational functions each require different techniques. Those techniques are the subject of Chapter 8.

\subsechead{Chapter summary}
This chapter built the integral from the ground up and bound it to the derivative. Section 6.1 met the area and distance problems and discovered them to be one limit of Riemann sums. Section 6.2 named that limit the \emph{definite integral} \(\int_{a}^{b} f(x)\,dx\) (Definition 6.2), established that a continuous integrand is integrable (Theorem 6.1, taken on credit), read the integral as net signed area, and collected its properties (Theorem 6.2).

Section 6.3 proved the \emph{Fundamental Theorem of Calculus}: Part 1 (Theorem 6.3) shows the accumulation function \(\int_{a}^{x} f\) is an antiderivative of \(f\), so differentiating an integral returns the integrand; Part 2 (Theorem 6.4) evaluates any definite integral from a single antiderivative. That the derivative of an accumulated area is the height, and that an area is therefore recovered from an antiderivative, is the central fact of the subject.

Section 6.4 introduced the notation \(\int f(x)\,dx\) for the \emph{indefinite integral} (Definition 6.4) and a working table, singled out \(\int \tfrac{1}{x}\,dx = \ln\lvert x\rvert + C\), and reread the theorem as the Net Change Theorem (Theorem 6.5). Section 6.5 turned the chain rule around into the Substitution Rule (Theorems 6.6–6.7) and used it to exploit symmetry (Theorem 6.8).

The results worth carrying forward are few and central:

\[ \int_{a}^{b} f(x)\,dx = \lim_{n \to \infty} \sum_{i=1}^{n} f(x_{i}^{*})\,\Delta x, \qquad \frac{d}{dx}\int_{a}^{x} f(t)\,dt = f(x), \]

\[ \int_{a}^{b} f(x)\,dx = F(b) - F(a), \qquad \int \frac{1}{x}\,dx = \ln\lvert x\rvert + C, \]

\[ \int f\bigl(g(x)\bigr)g'(x)\,dx = \int f(u)\,du. \]

The rest of the book uses the integral as a central tool. Chapter 7 applies the Fundamental Theorem to areas between curves, volumes, and arc lengths; Chapter 8 extends substitution into a full set of techniques for finding the antiderivatives that evaluation demands; and the integral returns, transformed, in the study of series, of differential equations, and of the calculus of several variables. Every one of those chapters rests on the two theorems proved in §6.3.
\end{document}
